%!TEX root = ../../main.tex
%% 1.5 인식론적 지위 — 근거: journal 02_METHODS_CORE §0, AGENTS.md, master_thesis_v1/README.md 집필 규칙
\section{증거의 지위와 집필 규칙}
\label{sec:intro-epistemic}

본 논문의 모든 수치는 동반 저널 원고와 같은 실험 원장에서 옮긴 것이다. 학위논문을 위해 새 실험, 재집계, 재학습을 수행하지 않았으며, 따라서 학위논문과 저널이 서로 다른 수치를 갖는 일은 없다. 저널이 예측 가능성과 검증을 동결된 측정 아래에서 압축해 보고한다면, 학위논문은 같은 증거를 설계 근거, 대안, 한계 중심으로 장 단위로 확장한다.

증거의 인식론적 지위도 그대로 밝힌다. 승률·시간 기준선을 유연하게 만든 주 대비, 점수 분해, 외부 이중 단계 평가는 초기 분석에서 시험 패치 15.16이 이미 노출된 뒤 심사 대응의 형태로 수행되었다. 따라서 해당 표들은 잠긴 가중치와 문서화된 절차로 질문에 답하지만, 사전 등록된 확증 검정이 아니라 \emph{탐색적} 결과로 표기한다(설계 태그 \texttt{EXPLORATORY\_REVIEW\_RESPONSE\_PRIOR\_TEST\_EXPOSURE}). 소규모 교전 코호트 S는 2026년 9월 21일 예비 선언된 계약 아래 추가되었고, 그 15.16 행 역시 이전 계보에서 점수화된 적이 있으므로 T와 같은 지위를 갖는다. 실행이 서면 설계와 달랐던 부분(균일 매듭, 평균 정규화하지 않은 가중치, 두 단계 선택)은 역사를 고쳐 쓰지 않고 실행 부록의 기록을 따른다.

집필에서는 계획만 있고 실행되지 않은 항목(정보군 모듈 M-F0--M-F3, 외부 확률 어댑터, 참여 인원 셀, 정의 상수의 전면 재민감도)을 ``본 연구에서 보였다''고 쓰지 않고, 제 \ref{ch:discussion} 장과 부록 \ref{app:deferred}에 미실행 계획으로만 둔다. 선행 연구의 교전 승자 예측 계보(\texttt{market\_event} 라벨)는 정의의 계보와 대응 검사에만 인용하며 본 논문의 주 결과로 쓰지 않는다.
